\documentclass[11pt]{article}
\usepackage[a4paper,margin=1in]{geometry}
\usepackage[T1]{fontenc}
\usepackage[utf8]{inputenc}
\usepackage{lmodern}
\usepackage{amsmath,amssymb,amsthm,mathtools}
\usepackage{booktabs}
\usepackage{microtype}
\usepackage[hidelinks]{hyperref}
\usepackage{enumitem}
\usepackage{array}

\newtheorem{theorem}{Theorem}
\newtheorem{lemma}{Lemma}
\newtheorem{proposition}{Proposition}
\newtheorem{corollary}{Corollary}
\theoremstyle{definition}
\newtheorem{definition}{Definition}
\theoremstyle{remark}
\newtheorem{remark}{Remark}

\newcommand{\Acal}{\mathcal A}
\newcommand{\tauac}{\tau_{\rm ac}}

\title{A Phase Transition on the Third Subset Diagonal\\of the Avoid-or-Compress Lemma}
\author{Ryutaro Yonezu\\Independent Researcher}
\date{22 September 2026\\Revised preprint -- not peer reviewed}

\begin{document}
\maketitle

\begin{abstract}
We determine the exact worst-case avoid-or-compress distance on the third subset diagonal $|S|-|A|=3$ for synchronizing strongly connected automata. If $s=|S|$, then the tight value is
\[
T_3(n,s)=
\begin{cases}
1,&s=n,\\
3,&s=n-1,\\
n-s+3,&s\le n-2.
\end{cases}
\]
The first two rows are forced by classical compression: a synchronizing automaton has a singular letter, and an $(n-1)$-subset can be compressed within three steps. For every $s\le n-2$ we construct an explicit binary, strongly connected, synchronizing automaton attaining the general avoid-or-compress upper bound $n-|A|=n-s+3$. The proper-subset constructions split according to $L=n-|A|$ modulo three. Their lower bound is governed by a three-token corridor on a cyclic block together with a constant-size tail defect. Synchronization is proved symbolically using full cycles and local merges, or an explicit reset word. An independent verifier checks 1,680 constructions with $1\le |A|\le30$ and $5\le n-|A|\le60$, up to $n=90$, and verifies the exact distance, strong connectivity, pair synchronizability, and the identities used in the synchronization proof. The result exhibits an exact one-step phase transition between the boundary $|S|=n-1$ and the proper region $|S|\le n-2$.
\end{abstract}

\section{Introduction}
The avoid-or-compress lemma is a useful local tool in synchronizing automata: for nonempty $A\subsetneq S\subseteq Q$, it bounds the length required either to remove at least one state of $A$ from the image of $S$ or to compress $S$ itself. In the synchronizing setting, the general upper bound is $n-|A|$ \cite{Szykula2018,Szykula2026}. The recent survey of open problems asks for parameter-dependent tight bounds and stresses that extremal examples should ideally remain synchronizing and strongly connected \cite{Szykula2026}.

The second subset diagonal $|S|-|A|=2$ admits exact tightness throughout the proper-subset region \cite{Yonezu2026Second}. The third diagonal is the first place where a different phenomenon appears. The classical Frankl--Pin one-step subset-compression bound
\[
\binom{n-|S|+2}{2}
\]
can be smaller than $n-|A|$. Exactly one boundary layer is affected: when $|S|=n-1$, compression gives $3$ whereas the avoid-or-compress lemma gives $4$. One step farther inside, at $|S|=n-2$, the two bounds both equal $5$, and the general avoid-or-compress value becomes attainable again.

We prove that this comparison is not merely an upper-bound artifact: it gives the exact phase transition. The lower-bound constructions use a common cyclic architecture. All states of $A$ form a chain under one letter; the complement $B$ is a cycle under the other. The three active $B$-states behave like tokens moving through a periodic coloring. Away from a constant-size defect near the end of the cycle, every use of the non-cyclic letter reconstructs the initial subset exactly. Thus success is delayed until the token window reaches the defect.

\section{Preliminaries}
Let $\Acal=(Q,\Sigma,\delta)$ be a DFA with $|Q|=n$. For $X\subseteq Q$ and $w\in\Sigma^*$, write $Xw$ for the image of $X$.

\begin{definition}
For $\varnothing\ne A\subsetneq S\subseteq Q$, define
\[
\tauac(\Acal;S,A)=\min\{|w|:A\nsubseteq Sw\text{ or }|Sw|<|S|\}.
\]
\end{definition}

On the third diagonal let
\[
|S|-|A|=3,
\]
and write
\[
r=|A|,\qquad L=n-r=n-|A|.
\]
Then $|S|=r+3$ and $n-|S|=L-3$.

Two general upper bounds are relevant. The avoid-or-compress lemma gives
\begin{equation}\label{eq:aoc}
\tauac(\Acal;S,A)\le L,
\end{equation}
while the Frankl--Pin subset-compression bound gives a compressing word of length at most
\begin{equation}\label{eq:fp}
\binom{n-|S|+2}{2}=\binom{L-1}{2}.
\end{equation}
The combinatorial extremal estimate underlying this compression theorem was proved by Frankl in response to a problem posed by Pin \cite{Frankl1982}; the resulting automata bound is commonly referred to as the Frankl--Pin or Pin--Frankl bound \cite{Volkov2022}. See also \cite{Szykula2018,Szykula2026}.

\section{Exact theorem and the boundary phase transition}
\begin{theorem}[Exact third diagonal]\label{thm:main}
Let $\Acal$ range over synchronizing strongly connected DFAs, and let $\varnothing\ne A\subsetneq S\subseteq Q$ satisfy $|S|-|A|=3$. Put $s=|S|$. The tight worst-case value is
\[
T_3(n,s)=
\begin{cases}
1,&s=n,\\[1mm]
3,&s=n-1,\\[1mm]
n-s+3,&s\le n-2.
\end{cases}
\]
Moreover, every lower bound in the last two rows is attained by a binary, strongly connected, synchronizing automaton.
\end{theorem}

\begin{proof}[Upper bounds]
If $s=n$, synchronization implies the existence of a singular letter, hence $T_3(n,n)\le1$; the empty word is not successful, so equality holds.

If $s=n-1$, \eqref{eq:fp} gives a compressing word of length at most $\binom32=3$, hence $T_3(n,n-1)\le3$.

If $s\le n-2$, the general avoid-or-compress bound gives $T_3(n,s)\le n-s+3$. The constructions below meet these bounds.
\end{proof}

\subsection{The $s=n-1$ lower bound}
Take the standard binary $n$-state \v{C}ern\'y automaton with states $q_0,\ldots,q_{n-1}$, where $a$ is the $n$-cycle and $b$ fixes all states except $q_{n-1}b=q_0$. Let
\[
S=Q\setminus\{q_{n-1}\},\qquad A=\{q_2,\ldots,q_{n-3}\}.
\]
For $n\ge5$, $|S|-|A|=3$. No word of length at most two either compresses $S$ or removes a state of $A$, whereas
\[
Sa^3=Q\setminus\{q_2\}
\]
avoids $q_2$. Thus $\tauac=3$. The automaton is binary, strongly connected, and synchronizing.

\section{Proper-subset constructions}
We now assume $s\le n-2$, equivalently $L\ge5$. Let
\[
A=\{A_0,\ldots,A_{r-1}\},\qquad B=\{B_0,\ldots,B_{L-1}\},
\]
\[
Q=A\cup B,\qquad S=A\cup\{B_0,B_1,B_2\}.
\]
Letter $b$ fixes $A$ and cycles $B$:
\[
A_i b=A_i,\qquad B_j b=B_{j+1\bmod L}.
\]
On $A$, letter $a$ is the chain
\[
A_0\to A_1\to\cdots\to A_{r-1}\to B_0.
\]
The action of $a$ on $B$ depends on $L\bmod3$.

\subsection{$L\equiv2\pmod3$}
For $0\le j<L-2$, set
\[
B_j a=\begin{cases}
B_1,&j\equiv0\pmod3,\\
A_0,&j\equiv1\pmod3,\\
B_2,&j\equiv2\pmod3,
\end{cases}
\]
and at the tail set
\[
B_{L-2}a=B_4,\qquad B_{L-1}a=A_0.
\]
A shortest successful word is
\begin{equation}\label{eq:w2}
w_2=b^{L-2}a^2.
\end{equation}

\subsection{$L\equiv0\pmod3$}
For $0\le j<L-1$, set
\[
B_j a=\begin{cases}
A_0,&j\equiv0\pmod3,\\
B_2,&j\equiv1\pmod3,\\
B_1,&j\equiv2\pmod3,
\end{cases}
\]
and set
\[
B_{L-1}a=B_{L-1}.
\]
A shortest successful word is
\begin{equation}\label{eq:w0}
w_0=b^{L-3}aba.
\end{equation}

\subsection{$L\equiv1\pmod3$}
Here $L\ge7$. For $0\le j<L-1$, set
\[
B_j a=\begin{cases}
B_1,&j\equiv0\pmod3,\\
B_2,&j\equiv1\pmod3,\\
A_0,&j\equiv2\pmod3,
\end{cases}
\]
and define the single tail defect
\[
B_{L-1}a=B_{L-3}.
\]
A shortest successful word is
\begin{equation}\label{eq:w1}
w_1=b^{L-2}a^2.
\end{equation}

\section{Lower bound: three-token corridor and defect zone}
The lower-bound proof is the same in spirit in all three residue classes. The $A$-chain contributes
\[
A\setminus\{A_0\}\quad\text{and}\quad B_0
\]
after applying $a$. Hence, for a safe subset of the form $A\cup T$ with $T\subseteq B$ and $|T|=3$, the three $B$-images under $a$ must provide $A_0$ together with two distinct $B$-states different from $B_0$. In the regular part of each construction, every three consecutive $B$-states do exactly this.

\begin{lemma}[Corridor reset]\label{lem:corridor}
Suppose the moving window
\[
W_k=\{B_k,B_{k+1},B_{k+2}\}
\]
does not intersect the residue-class tail defect. Then
\[
(A\cup W_k)a=S.
\]
Consequently, any occurrence of $a$ in the regular corridor erases the preceding displacement and returns the active subset exactly to $S$.
\end{lemma}
\begin{proof}
In each residue class, the regular period has length three and the three images are precisely a permutation of $\{A_0,B_1,B_2\}$. The $A$-chain supplies $A_1,\ldots,A_{r-1},B_0$. Their union is $S$.
\end{proof}

Thus a shortest successful word cannot use its first $a$ while the three-token window is in the regular corridor: deleting the prefix through that $a$ would produce a shorter successful word. The first $a$ of a shortest witness must occur only after the window reaches the final constant-size defect zone.

\begin{lemma}[Defect-zone check]\label{lem:defect}
For each of the three residue classes, every word of length less than $L$ whose first $a$ occurs in the defect zone is safe. The words \eqref{eq:w2}, \eqref{eq:w0}, and \eqref{eq:w1}, respectively, are successful and have length $L$.
\end{lemma}
\begin{proof}
For compactness write
\[
X(i,j,k)=A\cup\{B_i,B_j,B_k\},
\]
with $B$-indices interpreted modulo $L$. Every displayed $X(i,j,k)$ is safe because it contains all of $A$ and exactly three distinct $B$-states. Before the first occurrence of $a$ a word is necessarily a power $b^t$. By Lemma~\ref{lem:corridor}, every value of $t$ not listed below is in the regular corridor and the first $a$ returns the image to $S$. It therefore suffices to inspect the final constant number of values of $t$ and the suffixes that still fit inside a word of total length $<L$.

\smallskip
\noindent\emph{Case $L\equiv2\pmod3$.}
First assume $L\ge8$. The only defect displacements relevant below length $L$ are $t=L-4,L-3,L-2$. For $t=L-4$ the first $a$ gives $X(0,2,4)$ and at most two further letters are available. Direct substitution gives
\[
\begin{array}{c|ccccccc}
u&\varepsilon&a&b&aa&ab&ba&bb\\ \hline
X(0,2,4)u&X(0,2,4)&S&X(1,3,5)&S&X(1,2,3)&S&X(2,4,6).
\end{array}
\]
All seven images are safe. For $t=L-3$ the first image is again $X(0,2,4)$ but only suffixes of length at most one are possible, so this is already covered. For $t=L-2$ the first image is $X(0,1,4)$ and no further letter is allowed below length $L$. The remaining displacement $t=L-1$ places the first $a$ at length $L$ itself. The exceptional value $L=5$ has the same conclusion; wrap-around changes only
\[
X(0,2,4)b=X(0,1,3),\quad
X(0,2,4)ba=X(0,1,4),\quad
X(0,2,4)bb=X(1,2,4),
\]
and these images are still safe. Finally, from $X(0,1,4)$ a further $a$ compresses, so $b^{L-2}a^2$ succeeds at length $L$.

\smallskip
\noindent\emph{Case $L\equiv0\pmod3$.}
The only relevant defect displacements are $t=L-3$ and $t=L-2$. In both cases the first $a$ gives
\[
X(0,2,L-1).
\]
For $t=L-3$ there is room for at most one further letter, and
\[
X(0,2,L-1)a=X(0,1,L-1),\qquad
X(0,2,L-1)b=X(0,1,3),
\]
both safe. For $t=L-2$ no further letter is available below length $L$. The next displacement puts the first $a$ at length $L$. Moreover,
\[
X(0,2,L-1)b=X(0,1,3),
\]
and a final $a$ compresses this set, proving that $b^{L-3}aba$ succeeds at length $L$.

\smallskip
\noindent\emph{Case $L\equiv1\pmod3$.}
Here $L\ge7$. The only relevant defect displacements are $t=L-3$ and $t=L-2$. At $t=L-3$ the first $a$ gives
\[
X(0,2,L-3),
\]
and the only possible one-letter continuations below length $L$ satisfy
\[
X(0,2,L-3)a=S,\qquad
X(0,2,L-3)b=X(1,3,L-2),
\]
so both remain safe. At $t=L-2$ the first $a$ gives $X(0,1,L-3)$ and no further letter is available in a word shorter than $L$. A second $a$ maps two of its three $B$-states to the same state and compresses, so $b^{L-2}a^2$ succeeds at length $L$.

These lists exhaust every location of the first $a$ and every continuation compatible with total length $<L$. Hence every shorter word is safe in all three residue classes, while the displayed length-$L$ words succeed. The verifier in Section~\ref{sec:audit} independently reproduces the same finite transitions over the full parameter grid.
\end{proof}

\begin{proposition}[Exact proper-subset distance]\label{prop:exact}
For every $r\ge1$ and $L\ge5$, the corresponding construction satisfies
\[
\tauac(\Acal;S,A)=L=n-|A|.
\]
\end{proposition}
\begin{proof}
Lemma \ref{lem:corridor} forces a shortest successful word into the tail defect, and Lemma \ref{lem:defect} excludes success before length $L$ while providing a successful word of length exactly $L$.
\end{proof}

\section{Strong connectivity}
\begin{proposition}\label{prop:sc}
Every proper-subset construction is strongly connected.
\end{proposition}
\begin{proof}
Every $A_i$ reaches $B_0$ by powers of $a$, and $b$ traverses the entire $B$-cycle. In every residue-class construction at least one $B_j$ satisfies $B_j a=A_0$. Thus every state reaches $A_0$, and $A_0$ reaches every state.
\end{proof}

\section{Synchronization}
We use the following standard observation.

\begin{lemma}[Cycle-and-merge]\label{lem:cyclemerge}
Let $C$ be a finite set on which a word $\alpha$ acts as a full cycle. If the transformation semigroup contains an idempotent that fixes all but one point of $C$ and sends that point to an adjacent point on the $\alpha$-cycle, then $C$ is synchronizable. The same holds if an idempotent fixes all but the last two consecutive points and maps both of those points to their common predecessor.
\end{lemma}
\begin{proof}
Conjugating the local merge by powers of the full cycle moves it to every cyclic position. Successively applying these conjugates reduces the image size until one state remains.
\end{proof}

\subsection{$L\equiv1\pmod3$}
Let
\[
C=A\cup\{B_0,B_1,B_2\},\qquad m=r+3.
\]
Directly,
\[
Qa^2=C,
\]
and $a$ acts on $C$ as the full $m$-cycle
\[
A_0\to\cdots\to A_{r-1}\to B_0\to B_1\to B_2\to A_0.
\]
Set
\[
f=b^{L-2}a^2,\qquad \delta=f a^{m-2},\qquad e=\delta^2.
\]
A direct substitution shows that $e$ fixes every core position except the penultimate one, which it sends to the last. By Lemma \ref{lem:cyclemerge}, $C$ and hence $Q$ synchronize.

\subsection{$L\equiv2\pmod3$}
Again $Qa^2=C=A\cup\{B_0,B_1,B_2\}$. Moreover,
\[
C b^{L-2}a^2=D:=A\cup\{B_0,B_1\},
\]
and $a$ is a full cycle on $D$, of size $m=r+2$.

For $L\ge8$, set
\[
\gamma=b^2ab^{L-2}a^2,\qquad \delta=\gamma a^{m-3},\qquad e=\delta^2.
\]
Then $e$ fixes positions $0,\ldots,m-3$ and maps the last two positions to position $m-3$. Lemma \ref{lem:cyclemerge} synchronizes $D$. For the exceptional case $L=5$, the word $\gamma=b^2aba^2$ yields, after the same normalization, a standard adjacent idempotent merge.

\subsection{$L\equiv0\pmod3$}
Here an explicit reset word is available:
\begin{equation}\label{eq:reset0}
R=ab^{L-2}a\,(b^{L-1}a)^{r+1}ba.
\end{equation}
After the prefix $ab^{L-2}a$, the image is
\[
X_0=A\cup\{B_0,B_2,B_{L-1}\}.
\]
Let $T=b^{L-1}a$. For $0\le k<r$,
\[
X_k=\{A_k,\ldots,A_{r-1}\}\cup\{B_0,B_2,B_{L-1}\}
\]
satisfies $X_kT=X_{k+1}$. After $r+1$ applications of $T$, the image is $\{B_2,B_{L-1}\}$, and
\[
\{B_2,B_{L-1}\}ba=\{A_0\}.
\]
Hence \eqref{eq:reset0} is a reset word.

Combining these three cases with Proposition \ref{prop:sc} completes the lower-bound construction required for Theorem \ref{thm:main}.

\section{Computational audit}\label{sec:audit}
An independent verifier constructs every residue-class family and checks:
\begin{enumerate}[label=(\roman*)]
\item exact breadth-first first-success distance $L$;
\item success of the closed-form length-$L$ witness;
\item strong connectivity;
\item synchronizability via the exact unordered-pair criterion;
\item the symbolic core identities used in the synchronization proof.
\end{enumerate}
For the grid
\[
1\le r\le30,\qquad 5\le L\le60,
\]
there are 1,680 proper-subset constructions, up to $n=90$. All 1,680 pass every check. The boundary $L=3$ is the full-set case, and $L=4$ is handled by the compression upper bound and the \v{C}ern\'y witness.

\begin{center}
\begin{tabular}{lr}
\toprule
Audit item & Result\\
\midrule
Proper-subset parameter pairs & 1,680\\
Exact first-success distance & 1,680 / 1,680\\
Closed-form witness & 1,680 / 1,680\\
Strong connectivity & 1,680 / 1,680\\
Pair synchronizability & 1,680 / 1,680\\
Proof-level core identity & 1,680 / 1,680\\
Failures & 0\\
\bottomrule
\end{tabular}
\end{center}

\section{Discussion}
The third diagonal is the first diagonal on which the avoid-or-compress upper bound and the classical one-step compression bound cross. The exact theorem shows that the crossing produces a genuine extremal phase transition, not merely a proof artifact:
\[
1,\ 3,\ 5,\ 6,\ 7,\ldots
\]
are the successive worst-case values as the codimension $n-|S|$ grows from $0$ upward. In particular, the general value $n-|A|$ becomes sharp immediately once $|S|\le n-2$.

The next unresolved point is no longer the third diagonal itself but the interaction between diagonal index and codimension in the region where the Frankl--Pin bound competes closely with avoid-or-compress. The companion structural paper studies the concrete equality candidate $(d,k)=(3,7)$, where both general bounds equal ten.

\section*{Data and code availability}
The accompanying script \texttt{NO18\_THIRD\_DIAGONAL\_VERIFY\_V3.py} reproduces the 1,680-case audit and the symbolic synchronization identities.

\begin{thebibliography}{9}
\bibitem{Szykula2018}
M. Szyku\l a,
\newblock Improving the Upper Bound on the Length of the Shortest Reset Word,
\newblock \emph{STACS 2018}, LIPIcs 96, Article 56, 2018.
\newblock doi:10.4230/LIPIcs.STACS.2018.56.

\bibitem{Szykula2026}
M. Szyku\l a,
\newblock Synchronizing Automata: Open Problems,
\newblock \emph{Electronic Proceedings in Theoretical Computer Science} 451 (2026), 33--47.
\newblock doi:10.4204/EPTCS.451.3; arXiv:2608.24245.

\bibitem{Frankl1982}
P. Frankl,
\newblock An Extremal Problem for two Families of Sets,
\newblock \emph{European Journal of Combinatorics} 3(2) (1982), 125--127.
\newblock doi:10.1016/S0195-6698(82)80025-5.

\bibitem{Yonezu2026Second}
R. Yonezu,
\newblock Exact Tightness on the Second Subset Diagonal of the Avoid-or-Compress Lemma,
\newblock GitHub release v1.0.0, 2026.
\newblock \url{https://github.com/yonezaemon1-hub/second-diagonal-avoid-or-compress-tightness/releases/tag/v1.0.0}.
\newblock Archival DOI pending.

\bibitem{Volkov2022}
M. V. Volkov,
\newblock Synchronization of finite automata,
\newblock \emph{Russian Mathematical Surveys} 77(5) (2022), 819--891.

\bibitem{Ferens2021}
R. Ferens, M. Szyku\l a, and V. Vorel,
\newblock Lower Bounds on Avoiding Thresholds,
\newblock \emph{MFCS 2021}, LIPIcs 202, Article 46, 2021.
\newblock doi:10.4230/LIPIcs.MFCS.2021.46.
\end{thebibliography}

\end{document}
